\section{Quantization and Deployment}
\label{sec:quantization}

The deploy stack uses ternary weight packing plus INT8 activation quantization for efficient artifacts.

\begin{figure}[t]
\centering
\fitfigure{\begin{tikzpicture}[node distance=1.3cm and 0.9cm]
\node[box, fill=blue!8] (ckpt) {Trained checkpoint};
\node[box, fill=yellow!10, right=of ckpt] (quant) {Weight packing\\ternary / low-bit};
\node[box, fill=orange!12, right=of quant] (act) {Activation scaling\\INT8 path};
\node[box, fill=green!8, right=of act] (artifact) {Deployable artifact\\smaller storage footprint};
\draw[line] (ckpt) -- (quant);
\draw[line] (quant) -- (act);
\draw[line] (act) -- (artifact);
\end{tikzpicture}}
\caption{Deployment path from a trained checkpoint to a compact artifact. The codebase treats quantization as a deploy-stage transformation rather than a separate model family.}
\label{fig:deploy-flow}
\end{figure}

\subsection{Ternary Quantization and Activation Scaling}
The deploy-stage transforms a trained checkpoint into a compact artifact via BitNet-style ternary weight packing \citep{wang_bitnet_2023}, INT8 activation scaling, and 1.58-bit storage approximation aligned with lightweight deployment goals \citep{samson_lightweight_2026,bravin_embbert_2026}. The mathematical formulations of the ternary weight scaling, the INT8 activation quantization/dequantization, and the FP32/FP16/1.58-bit size estimates are deferred to Appendix~\ref{sec:annex-norm-bitnet-mod}.